\documentclass[12pt]{article}

\usepackage[T1]{fontenc}
\usepackage[utf8]{inputenc}
\usepackage[brazil]{babel}
\usepackage{lmodern}
\usepackage{microtype}
\usepackage[a4paper,margin=2.5cm]{geometry}
\usepackage{setspace}
\usepackage{indentfirst}

\setstretch{1.15}

\title{A Síntese Novo-Keynesiana \thanks{Texto Traduzido livremente por Luiz Mario Andrade com o auxílio da ferramenta Chat GPT 5.4}}
\author{David Romer}
\date{Inverno de 1993}

\begin{document}

\maketitle

Duas crenças sobre a economia impulsionaram o afastamento da macroeconomia keynesiana da ortodoxia clássica na década de 1930. A primeira era de que havia um desemprego involuntário generalizado: muitas pessoas pareciam dispostas a trabalhar, mas incapazes de encontrar emprego ao salário vigente. A segunda era de que as flutuações na demanda agregada eram uma fonte central de mudanças de curto prazo na atividade econômica agregada: mudanças na demanda do governo por bens, na confiança dos líderes empresariais e nos mercados monetários e financeiros pareciam ter efeitos poderosos sobre o emprego e o produto. Uma nova teoria era necessária para incorporar essas crenças.

A famosa "síntese neoclássica", que se desenvolveu ao longo das três décadas seguintes, postulou uma única explicação para ambos os fenômenos: que os preços em unidades monetárias se ajustavam apenas lentamente aos desequilíbrios entre oferta e demanda. O mais importante desses preços monetários rígidos era o preço monetário do trabalho — o salário nominal. O ajuste salarial lento implicava que a demanda e a oferta de trabalho poderiam estar em desequilíbrio e, portanto, que o desemprego poderia surgir. Além disso, como eram os salários nominais que demoravam a se ajustar, a nova teoria implicava que a dicotomia clássica entre variáveis nominais e reais falhava e, portanto, que movimentos em variáveis nominais, como a oferta de moeda, poderiam ter grandes efeitos sobre variáveis reais, como o produto e o emprego.

O restante da síntese neoclássica era walrasiano. Os mercados de bens e trabalho eram competitivos, as externalidades estavam ausentes e a informação era perfeita. A síntese atingiu seu auge com os modelos de desequilíbrio do final dos anos 1960 e início dos anos 1970, que anexaram suposições de preços e salários completamente fixos a modelos de equilíbrio geral walrasianos (por exemplo, Malinvaud, 1977).

A síntese neoclássica naufragou no que parece, em retrospecto, uma pergunta óbvia: em um ambiente que é tão implacavelmente competitivo, como pode persistir um desvio tão flagrante do comportamento walrasiano? Talvez a mensagem mais fundamental da economia seja a de que, em um ambiente competitivo, incentivos poderosos impelem os atores econômicos a ajustar os preços em resposta a desequilíbrios entre a oferta e a demanda.

Essa questão levou ao colapso da síntese neoclássica e gradualmente dividiu a macroeconomia convencional em duas escolas. Uma escola — a teoria dos ciclos reais de negócios — abandonou não apenas a síntese neoclássica, mas as premissas da macroeconomia keynesiana. Os membros desta escola negam a existência de desemprego involuntário significativo e de quaisquer falhas importantes da dicotomia clássica (ver, por exemplo, Plosser, 1989).

A outra escola — a macroeconomia novo-keynesiana — também passou a reconhecer que a microeconomia walrasiana e a síntese neoclássica não forneciam fundamentos teóricos adequados para a macroeconomia keynesiana. Mas os membros desta escola acreditavam que a resposta apropriada era tentar determinar se uma descrição correta da microeconomia daria origem aos fenômenos que eles acreditavam caracterizar a macroeconomia. Afinal, aceitar a crença de que o mercado de trabalho estava continuamente em equilíbrio walrasiano exigiria negar que o desemprego fosse um fenômeno importante. E aceitar a dicotomia clássica exigiria negar que distúrbios monetários tivessem efeitos reais. Também exigiria desistir de uma explicação direta de como outras mudanças no lado da demanda — como mudanças nos gastos do governo e na demanda de investimento — poderiam ter efeitos reais substanciais. Na ausência de um ajuste lento de preços nominais, tais mudanças de demanda afetam a atividade real agregada apenas via o impacto da taxa de juros real e da riqueza sobre a oferta de trabalho, e a ideia de que esses efeitos eram grandes parecia implausível.

Os novos keynesianos progrediram mais rapidamente na compreensão da microeconomia do desemprego do que na compreensão da microeconomia da rigidez de preços nominais. Mas os últimos cinco anos viram avanços importantes nesta segunda área. Este artigo descreverá esses avanços, discutirá nosso entendimento atual da rigidez nominal e avaliará o trabalho que ainda resta a ser feito.

\section*{Fricções Nominais}

O elemento central da síntese neoclássica era sua suposição de que os preços não se ajustavam imediatamente para equilibrar a oferta e a demanda. A resposta natural ao colapso da síntese foi, portanto, investigar se o ajuste imperfeito de preços poderia ser derivado de suposições realistas sobre o ambiente microeconômico, em vez de ser apenas assumido.

A pesquisa resultante levou a uma variedade de teorias não walrasianas sobre o funcionamento dos mercados. Análises não walrasianas do mercado de trabalho, por exemplo, sugeriram que os salários poderiam servir a uma variedade de funções além de equilibrar a oferta e a demanda. Em modelos iniciais de contratos implícitos, o salário serve como um meio para a firma fornecer seguro aos seus trabalhadores; em muitos modelos de barganha, o salário é o meio pelo qual as rendas são divididas entre os trabalhadores e a firma; e em modelos de salário de eficiência, o salário afeta a produtividade do trabalho.

Modelos como esses têm o potencial de fornecer uma explicação para o desemprego. Mas eles não podem fornecer uma explicação para as falhas da dicotomia clássica. Os modelos focam em imperfeições reais: os trabalhadores estão preocupados em assegurar seus padrões de vida reais; firmas e trabalhadores estão preocupados com as rendas reais que obtêm; a produtividade do trabalho depende do salário real que a firma paga; e assim por diante. Se a oferta de moeda muda nesses modelos, então, assim como em modelos completamente clássicos, todos os preços nominais mudam, deixando os preços relativos e os resultados reais (com quaisquer características não walrasianas que possam envolver) inalterados.

Qualquer base microeconômica para a falha da dicotomia clássica requer algum tipo de imperfeição nominal; caso contrário, um distúrbio puramente nominal deixa o equilíbrio real (ou o conjunto de equilíbrios reais) inalterado.\footnote{Em suas contribuições a este simpósio, James Tobin e Bruce Greenwald e Joseph Stiglitz argumentam que, na margem, aumentar a velocidade do ajuste nominal pode tornar os efeitos reais dos distúrbios nominais maiores. Mesmo que isso esteja correto, a flexibilidade nominal completa e imediata tornaria os choques monetários e outros distúrbios da demanda agregada neutros. Assim, a compreensão das fontes da flexibilidade nominal incompleta permaneceria uma parte necessária de um relato das flutuações macroeconômicas.} Isso levanta imediatamente uma dificuldade. Os indivíduos estão, em última análise, preocupados com preços e quantidades reais: salários reais, horas de trabalho, níveis de consumo real e afins. As magnitudes nominais importam para eles apenas de maneiras que são menores e facilmente superáveis. Preços e salários são cotados em termos nominais, mas custa pouco mudá-los (ou indexá-los). Os indivíduos não são totalmente informados sobre o nível de preços agregado ou a oferta de moeda, mas podem obter informações bastante precisas com pouco custo. Contratos de dívida são geralmente especificados em termos nominais, mas eles também poderiam ser indexados com pouca dificuldade. E os indivíduos detêm quantias modestas de moeda, que é denominada em termos nominais, mas podem mudar suas participações facilmente. De modo algum as magnitudes nominais são de grande importância direta para os indivíduos. De fato, a dificuldade é revelada pela própria palavra: chamar algo de "nominal" é dizer que é meramente um nome.

Assim, se a falha da dicotomia clássica é importante para as flutuações na atividade agregada, deve ser que fricções nominais que parecem pequenas ao nível das famílias e firmas individuais — como o fato de que os preços são postados em unidades nominais, ou que obter informações precisas sobre o preço agregado envolve um custo — tenham de alguma forma um grande efeito na macroeconomia. É esse insight, devido a Mankiw (1985) e Akerlof e Yellen (1985), que levou ao progresso recente na compreensão das fundações microeconômicas do impacto real dos distúrbios da demanda agregada.\footnote{A natureza do insight também ajuda a explicar por que o progresso na compreensão da não-neutralidade dos choques de demanda agregada foi tão lento: é natural começar a busca por uma explicação de um grande enigma explorando grandes desvios da ortodoxia prevalecente, em vez de procurar grandes efeitos a partir de pequenos desvios.}

\section*{Rigidezes Reais}

A questão de saber se pequenas fricções podem fazer com que distúrbios nominais tenham grandes efeitos na atividade econômica agregada depende dos incentivos das firmas individuais para mudar seus preços quando o produto agregado muda. Como exemplo, considere um declínio do produto em toda a economia. A questão que se apresenta a uma firma quando a demanda por seu produto cai como resultado do declínio no produto agregado é se deve manter seu preço fixo e reduzir a produção, ou baixar seu preço e, assim, reduzir ou eliminar a necessidade de reduzir a produção.

Esta questão pode ser analisada usando o diagrama de receita marginal-custo marginal na Figura 1. A economia começa em equilíbrio; assim, a firma representativa está produzindo no ponto onde o custo marginal é igual à receita marginal (Ponto A no diagrama). Uma contração do produto em toda a economia desloca a curva de demanda que a firma enfrenta para dentro — a um dado preço, a demanda pelo produto da firma é menor. Assim, a curva de receita marginal se desloca para dentro. Se a firma não alterar seu preço, sua produção é determinada pela demanda ao preço existente (Ponto B). Nesse nível de produção, a receita marginal excede o custo marginal e, portanto, a firma tem algum incentivo para baixar seu preço e aumentar a produção. Se a firma alterar seu preço, produzirá no ponto onde o custo marginal e a receita marginal são iguais (Ponto C). A área do triângulo sombreado no diagrama mostra os lucros adicionais a serem obtidos com a redução do preço e o aumento da quantidade produzida. Para que a firma esteja disposta a manter seu preço fixo, a área do triângulo deve ser pequena.

O diagrama revela um ponto crucial: o incentivo da firma para reduzir seu preço pode ser pequeno mesmo que ela seja muito prejudicada pela queda na demanda. A firma preferiria enfrentar a curva de demanda original, mais alta, mas, é claro, ela só pode escolher um ponto na nova curva de demanda. A firma pode descobrir que os ganhos de reduzir seu preço são pequenos, mesmo que o deslocamento em sua curva de demanda seja grande.

Se os ganhos para a firma ao cortar seu preço forem de fato pequenos, o comportamento de muitas dessas firmas enfrentando pequenas fricções no ajuste de preços pode fazer com que um distúrbio da demanda agregada tenha grandes efeitos reais. Suponha que o distúrbio subjacente seja um declínio na oferta de moeda ou algum outro deslocamento adverso da demanda agregada, e suponha provisoriamente que as firmas não cortem seus preços em resposta a este distúrbio. Nesta situação, o produto real agregado cai. Assim, a situação enfrentada pela firma representativa é como a descrita na figura. Se o incentivo da firma representativa para ajustar seu preço for pequeno e houver fricções no ajuste de preços, então o comportamento conjeturado das firmas de manter seus preços fixos é de fato um equilíbrio. Se, por outro lado, o incentivo para o ajuste de preços for grande, todas as firmas cortam seus preços; o resultado final é que o choque negativo da demanda agregada resulta apenas em preços mais baixos.

\begin{center}
\textbf{Figura 1} \\
\textbf{O Incentivo de uma Firma Representativa para Mudar seu Preço em Resposta a uma Mudança no Produto Agregado} \\
\textit{[Descrição do Diagrama: Gráfico de Preço vs. Quantidade mostrando curvas de demanda D e D', receita marginal MR e MR', e custo marginal MC. O ponto A é o equilíbrio inicial. O ponto B é o nível de produção se o preço for mantido após um choque. O ponto C é o novo ótimo se o preço for ajustado. Um triângulo sombreado entre B, C e a curva MC representa o lucro perdido por não ajustar o preço.]}
\end{center}

O incentivo de uma firma para mudar seu preço em resposta à queda na demanda — o tamanho do triângulo na Figura 1 — é determinado pelas respostas do custo marginal e da receita marginal à retração na demanda agregada. Considere primeiro o custo marginal. Como menos produto está sendo produzido, menos trabalho é demandado. Com uma curva de oferta de trabalho inclinada para cima, isso implica um declínio no salário real e, portanto, no custo marginal.\footnote{Se o trabalho for móvel no curto prazo, o declínio no custo marginal assume a forma de um deslocamento para baixo da curva de custo marginal causado por uma queda no salário em toda a economia; se o trabalho for imóvel, o declínio assume a forma de um movimento ao longo de uma curva de custo inclinada para cima. Por simplicidade, não desloquei a curva na figura.} O comportamento cíclico do custo marginal também depende do grau de retornos decrescentes de curto prazo ao trabalho; se o produto marginal do trabalho aumentar rapidamente à medida que o insumo trabalho diminui, a curva de custo marginal é íngreme mesmo que o salário real seja constante. Quanto mais o custo marginal cai quando a produção diminui, maior o incentivo da firma para baixar seu preço.

Considere agora a receita marginal. Quanto mais a curva de receita marginal se desloca para a esquerda, menor o incentivo da firma para baixar seu preço. O tamanho do deslocamento da curva de receita marginal depende do comportamento cíclico da elasticidade da demanda. Na figura, assume-se que a elasticidade da demanda que a firma enfrenta ao seu preço existente não muda quando o produto agregado muda. Neste caso, a receita marginal ao preço existente (que agora corresponde a um nível mais baixo de produção) não é afetada pela mudança no produto em toda a economia. Se a elasticidade da demanda ao preço existente cai quando o produto agregado diminui, o deslocamento na receita marginal é maior; se a elasticidade aumenta, o deslocamento é menor.

O arcabouço estabelecido na Figura 1 pode ser usado para demonstrar que simplesmente adicionar competição imperfeita e pequenas barreiras ao ajuste de preços à visão de mundo dominante das décadas de 1950 e 1960 não é suficiente para fornecer uma base microeconômica para a visão de que choques de demanda agregada são centrais para as flutuações econômicas.\footnote{O argumento de Mankiw e de Akerlof e Yellen de que pequenas fricções podem gerar rigidez nominal significativa é frequentemente interpretado como segue. Os lucros de uma firma imperfeitamente competitiva são uma função suave do preço que ela cobra. Assim, os lucros perdidos ao se afastar do preço ótimo são de segunda ordem no tamanho do desvio, e assim o custo de não mudar o preço em resposta a um deslocamento na demanda agregada é de segunda ordem no tamanho do deslocamento — geometricamente, a região sombreada na Figura 1 é um triângulo. Assim, conclui a interpretação, a competição imperfeita sozinha é suficiente para explicar como fricções "pequenas" no ajuste de preços são suficientes para fazer com que os preços permaneçam fixos diante de movimentos na demanda agregada. Este argumento (que nem Mankiw nem Akerlof e Yellen fazem) repousa em uma confusão de dois usos do termo "pequeno". O que é necessário para fornecer uma base microeconômica para a visão de que movimentos na demanda agregada são importantes para as flutuações macroeconômicas é uma demonstração de que fricções que são "pequenas" no sentido de representarem barreiras empiricamente plausíveis à flexibilidade nominal de preços são suficientes para evitar que os preços se ajustem totalmente em resposta a movimentos na demanda agregada do tamanho tipicamente observado em flutuações cíclicas. O fato de as fricções necessárias serem "pequenas" no sentido de serem de segunda ordem no tamanho dos movimentos da demanda agregada é simplesmente irrelevante para essa questão (Reaume, 1991).} A fonte da dificuldade reside no mercado de trabalho. Se a oferta de trabalho for relativamente inelástica — certamente a visão predominante há 20 anos, e provavelmente a visão predominante hoje — e se não houver desvios das suposições walrasianas além da presença de pequenas barreiras ao ajuste nominal, então o declínio no insumo trabalho associado ao declínio na produção leva a uma grande queda no salário real.

Neste caso, o custo marginal cai consideravelmente em recessões. Como resultado, a menos que a elasticidade da demanda também caia bruscamente, os incentivos das firmas para reduzir os preços são grandes. Cálculos rápidos para um modelo simples em que a competição imperfeita é o único desvio das suposições walrasianas mostram que, se a oferta de trabalho for relativamente inelástica, os incentivos das firmas para mudar seus preços diante de movimentos da demanda agregada de alguns pontos percentuais superam quaisquer barreiras plausíveis ao ajuste nominal (Ball e Romer, 1990).

Assim, para que a dicotomia clássica falhe, deve ser que o custo marginal não caia drasticamente em resposta a uma contração da produção impulsionada pela demanda, ou que a receita marginal caia bruscamente, ou alguma combinação dos dois. Em um nível mais geral, o incentivo para mudar o preço em resposta a uma mudança no produto em toda a economia pode ser expresso como uma função de dois fatores: o impacto da mudança no preço real que maximiza o lucro da firma e o custo para a firma de um determinado desvio de seu preço real em relação ao nível que maximiza o lucro. Para que o incentivo ao ajuste diante de flutuações impulsionadas pela demanda seja pequeno, ou os preços reais que maximizam o lucro devem responder pouco às mudanças no produto agregado — na terminologia de Ball e Romer (1990), o grau de "rigidez real" deve ser alto — ou desvios consideráveis dos preços que maximizam o lucro devem ter apenas custos pequenos. No modelo simples discutido acima, as grandes mudanças nos salários reais em resposta aos movimentos do produto agregado fazem com que os preços que maximizam o lucro sejam muito sensíveis ao produto — isto é, a rigidez real é baixa — e, portanto, o incentivo ao ajuste é grande. Tanto uma menor sensibilidade cíclica do custo marginal quanto uma maior sensibilidade cíclica da receita marginal aumentam a rigidez real e, assim, reduzem os incentivos das firmas para ajustar seus preços. Em suma, um modelo completo de grandes efeitos reais de distúrbios nominais requer tanto fricções nominais quanto rigidezes reais.

\section*{Fontes Potenciais de Rigidez Real}

O conhecimento econômico não progrediu ao ponto de termos uma visão clara de quais são as rigidezes reais mais importantes, mas pesquisas recentes sugeriram alguns candidatos. Em vez de examinar todos eles, descreverei brevemente quatro dos mais promissores. Todos eles têm implicações potencialmente importantes para assuntos muito além dos incentivos que as firmas têm para mudar seus preços em resposta a choques de demanda agregada, e todos são áreas ativas de pesquisa e debate.

A primeira área de pesquisa diz respeito a economias externas de escala decorrentes de "externalidades de mercado denso" (por exemplo, Diamond, 1982). Este trabalho investiga mecanismos através dos quais a compra de insumos e a venda de produtos finais podem ser mais fáceis em tempos de alta atividade econômica, quando o comércio está ativo e os mercados estão funcionando bem, do que em tempos de baixa atividade econômica. Esses efeitos agem para deslocar a curva de custo marginal para baixo em expansões e para cima em recessões.

Uma segunda linha de trabalho considera as imperfeições do mercado de capital decorrentes da informação imperfeita. Esses modelos começam observando que a assimetria de informações entre credores e devedores é um obstáculo apenas na busca de financiamento externo. Segue-se que, em uma situação de informação assimétrica, o financiamento interno é menos dispendioso do que o financiamento externo. Como as firmas têm lucros maiores e, portanto, mais fundos disponíveis para financiamento interno em períodos de expansão do que em recessões, as imperfeições do mercado de capitais tendem a tornar o custo do capital contracíclico; e como os custos de capital são um componente importante dos custos totais, isso age para fazer a curva de custo se mover em uma direção contracíclica (por exemplo, Bernanke e Gertler, 1989).

Uma terceira área de pesquisa foca no comportamento cíclico das elasticidades da demanda nos mercados de bens. Há uma variedade de razões pelas quais a elasticidade da demanda pode variar em resposta aos movimentos do produto agregado. Por exemplo, quando o produto agregado é alto, os efeitos de "mercado denso" podem tornar mais fácil para as firmas disseminar informações e para os consumidores adquiri-las. Isso poderia agir para tornar a elasticidade da demanda, e portanto a curva de receita marginal, mais procíclica, e assim reduziria os incentivos das firmas para ajustar seus preços em resposta aos movimentos da demanda agregada.

Nenhuma dessas três áreas de pesquisa diz respeito ao mercado de trabalho. Mas as rigidezes reais no mercado de trabalho parecem ser uma parte necessária da explicação dos efeitos reais dos distúrbios nominais. Como explicado na seção anterior, se o mercado de trabalho fosse walrasiano e a oferta de trabalho inelástica, os salários reais seriam altamente procíclicos. Se esse padrão se mantivesse na prática, as rigidezes reais em outras partes da economia (como as que acabei de discutir) teriam de ser extremamente fortes para superar o grande incentivo ao ajuste criado pelos salários fortemente procíclicos. No entanto, embora os analistas disputem o comportamento cíclico preciso dos salários reais, não há evidências de que eles sejam fortemente procíclicos. A quarta e mais importante área de pesquisa sobre rigidezes reais busca explicar essa observação.

Em um nível geral, os salários reais podem não ser altamente procíclicos por duas razões. Primeiro, a oferta agregada de trabalho de curto prazo poderia ser relativamente elástica; tal alta elasticidade de curto prazo poderia surgir, por exemplo, da substituição intertemporal e de não convexidades na oferta de trabalho. No entanto, tais modelos encontraram pouco apoio empírico (por exemplo, Altonji, 1986). Segundo, talvez algum tipo de imperfeição no mercado de trabalho faça com que os trabalhadores estejam fora de suas curvas de oferta de trabalho durante pelo menos parte do ciclo de negócios. Por exemplo, os modelos de salário de eficiência implicam que as firmas estabelecem salários reais acima dos níveis de equilíbrio de mercado. Esses modelos quebram, assim, o vínculo estreito entre a elasticidade da oferta de trabalho e a resposta dos salários reais aos distúrbios da demanda e, portanto, implicam que os salários reais podem não ser altamente procíclicos mesmo que a oferta de trabalho seja bastante inelástica. Outras imperfeições do mercado de trabalho, como a informação imperfeita e o monopólio bilateral decorrente da heterogeneidade entre trabalhadores e empregos, poderiam ter implicações semelhantes para os movimentos do salário real. Se imperfeições reais como essas fazem com que os salários reais respondam pouco aos distúrbios da demanda, elas reduzem consideravelmente os incentivos das firmas para variar seus preços em resposta a esses deslocamentos de demanda.

Além disso, a possibilidade de rigidezes reais substanciais no mercado de trabalho sugere que o canal através do qual pequenas barreiras ao ajuste nominal fazem com que distúrbios nominais tenham efeitos reais substanciais pode envolver a rigidez de salários nominais, em vez de preços nominais. Se os salários apresentarem rigidez real substancial, uma expansão impulsionada pela demanda leva a apenas pequenos aumentos nos salários reais ótimos. Como resultado, assim como pequenas fricções no ajuste de preços podem levar a uma rigidez substancial de preços nominais, pequenas fricções no ajuste de salários nominais podem levar a uma rigidez substancial de salários nominais.

\section{*{Bem-estar}

Em uma economia puramente walrasiana, o nível de produto que prevalece sob flexibilidade total de preços é o ótimo. Nesse arcabouço, qualquer desvio do produto normal — seja um boom ou uma recessão — reduz o bem-estar. Além disso, como os retornos privados de uma ação são iguais aos benefícios sociais da ação, o fato de as barreiras ao ajuste nominal de preços serem pequenas implica imediatamente que o custo de bem-estar de um distúrbio nominal deve ser pequeno: se o custo de bem-estar excedesse o custo de ajustar os preços, o retorno privado de mudar os preços e, assim, evitar que o distúrbio tivesse efeitos reais, excederia o custo, e então os preços seriam alterados.

Os modelos novo-keynesianos têm implicações de bem-estar muito diferentes. Esses modelos implicam uma assimetria entre booms impulsionados pela demanda e recessões impulsionadas pela demanda, com os booms aumentando o bem-estar e as recessões diminuindo-o. Além disso, sugerem a possibilidade de que fricções nominais possam levar a um nível de volatilidade ineficientemente alto e que a política de estabilização governamental possa, portanto, ser desejável.

A competição imperfeita sozinha é suficiente para implicar assimetria entre booms e recessões. Sob competição imperfeita, como o custo marginal está abaixo do preço, o nível de emprego que maximiza o lucro é ineficientemente baixo, e o nível de preço que maximiza o lucro é alto demais. Um aumento no produto a partir de seu nível de equilíbrio, em vez de reduzir o bem-estar, aproxima a economia do ótimo social. Um declínio no produto, em contraste, afasta a economia ainda mais do ótimo (Mankiw, 1985).

Além disso, a ineficiência que surge sob competição imperfeita implica que os preços podem permanecer fixos diante de um declínio na demanda agregada, mesmo que os custos da recessão resultante sejam muito maiores do que os custos de ajustar os preços. Com a competição imperfeita, as decisões de precificação das firmas têm externalidades. A externalidade pode ser pensada como operando através da demanda agregada. Uma decisão de aumentar o preço, por exemplo, através de seu impacto no nível de preço agregado, move a economia para cima na curva de demanda agregada e, assim, reduz o produto agregado; esse declínio no produto prejudica outras firmas ao reduzir a demanda por seus bens. Da mesma forma, cortes de preços por todas as firmas em resposta a uma queda na demanda agregada impediriam a queda do produto agregado. Mas, como descrito acima, o ganho da firma individual ao cortar seu próprio preço pode ser pequeno. Assim, é possível que a queda na demanda agregada leve a preços inalterados e a uma recessão dispendiosa, embora as barreiras ao ajuste de preços sejam pequenas.

Outras teorias podem levar a resultados semelhantes. Por exemplo, se os salários de eficiência causam desemprego, segue-se imediatamente que o produto marginal do trabalho excede o valor marginal do lazer. Assim, novamente, o nível de equilíbrio do produto é menor que o nível ótimo, booms e recessões têm efeitos assimétricos sobre o bem-estar, e pequenas barreiras ao ajuste de salários ou preços podem ser suficientes para fazer com que declínios na demanda agregada resultem em recessões dispendiosas.

Assim como as decisões das firmas que afetam o nível de produção têm externalidades, geralmente também existem externalidades decorrentes de suas decisões que afetam a volatilidade do produto. Especificamente, uma política por parte de uma firma de manter seu preço fixo em resposta a aumentos e quedas na demanda agregada aumenta a magnitude das flutuações agregadas; isso normalmente prejudica outros na economia. Assim, as flutuações agregadas são ineficientemente grandes, e a política de estabilização governamental tem o potencial de corrigir uma falha de mercado e aumentar o bem-estar (Ball e Romer, 1989).

Se os benefícios de bem-estar ao corrigir tais flutuações ineficientes são grandes, é uma questão em aberto. É possível que flutuações de produto impulsionadas pela demanda causem flutuações substanciais no bem-estar, mas pouca mudança no bem-estar médio. Mas os benefícios da política de estabilização poderiam ser grandes se declínios no produto em relação ao seu nível de equilíbrio tivessem grandes custos de bem-estar, mas aumentos tivessem apenas pequenos benefícios. Isso poderia ocorrer, por exemplo, se a desutilidade do trabalho fosse uma função fortemente crescente da quantidade de trabalho ofertada; nesse caso, os custos de utilidade do aumento da oferta de trabalho em um boom compensariam amplamente os benefícios do aumento da produção, enquanto os ganhos do aumento do lazer em uma recessão seriam pequenos em relação aos custos da produção perdida. Uma segunda possibilidade é que a redução da volatilidade da demanda agregada aumente o produto médio. Isso poderia ocorrer se a variabilidade agregada tivesse um impacto importante na escolha de capacidade das firmas e, consequentemente, no produto potencial (Meltzer, 1988). Outro canal através do qual a redução da volatilidade poderia aumentar o produto médio é sugerido pela visão tradicional de que choques negativos de demanda agregada reduzem principalmente o produto, enquanto choques positivos aumentam principalmente os preços. Se essa visão se provar correta, ela poderia implicar ganhos significativos de bem-estar a partir de uma política de estabilização contracíclica.

\section*{“Falha de Coordenação” e “Histerese”}

Minha discussão anterior sobre os incentivos das firmas para ajustar os preços em resposta aos movimentos da demanda agregada suprimiu o que parece ser uma possibilidade incômoda: as rigidezes reais podem ser tão fortes que o incentivo para baixar o preço em resposta a uma contração da produção em toda a economia não seja apenas pequeno, mas inexistente. Ou seja, quando a produção se contrai, a interseção das novas curvas de receita marginal e custo marginal pode ocorrer em um nível de produção inferior à quantidade agora demandada ao preço antigo. O incentivo da firma representativa seria então reagir ao choque negativo da demanda agregada aumentando seu preço e reduzindo ainda mais sua produção.

Nesta situação, existe mais de um nível "normal" de produto. Considere novamente a Figura 1. Na posição inicial, a economia está em equilíbrio: a firma representativa está produzindo a quantidade que produziria se pudesse definir seu preço livremente. Se uma queda na oferta de moeda, com os preços de outras firmas inalterados, cria um incentivo para a firma representativa aumentar seu preço e cortar ainda mais a produção, isso significa que a produção ótima da firma está mudando mais do que proporcionalmente em relação à quantidade média produzida por outras firmas. Contanto que a produção ótima seja uma função contínua da produção média de outras firmas e haja algum limite para a capacidade de produção das firmas, deve haver pelo menos dois outros níveis de produção de equilíbrio. Isso é mostrado na Figura 2.

\begin{center}
\textbf{Figura 2} \\
\textbf{Falha de Coordenação} \\
\textit{[Descrição do Diagrama: Gráfico mostrando a função de reação $y^*_i(y)$ cruzando uma linha de 45 graus em três pontos: A (equilíbrio inicial), B (equilíbrio superior) e C (equilíbrio inferior). O eixo vertical é $y^*_i$ (produção ótima da firma) e o eixo horizontal é $y$ (produção média das outras firmas).]}
\end{center}

A Figura 2 plota a produção ótima de uma firma representativa ($y^*_i$) como uma função da quantidade média produzida por outras firmas ($y$). Em qualquer ponto onde $y^*_i = y$ — a função de reação cruza a linha de 45 graus — a firma representativa deseja produzir a quantidade média que outras firmas estão produzindo e, portanto, a economia está em equilíbrio. O ponto A mostra a posição inicial: a economia está em equilíbrio e a produção ótima da firma representativa move-se mais do que proporcionalmente à produção média. Assumindo algum limite para a capacidade de produção das firmas, a função de reação deve cruzar a linha de 45 graus novamente em algum nível superior de produção (Ponto B). E, como a produção não pode ser negativa, deve haver uma terceira interseção em algum ponto abaixo de A (Ponto C). Além disso, se houver externalidades positivas de uma produção mais alta, como sob competição imperfeita, podemos fazer classificações de bem-estar dos vários equilíbrios: os equilíbrios com maior produção são superiores.

Modelos com múltiplos equilíbrios classificados por bem-estar são conhecidos como modelos de falha de coordenação (Cooper e John, 1988). Eles têm a capacidade potencial de fornecer uma explicação para as flutuações econômicas sem qualquer dependência de barreiras ao ajuste de preços, explicar "equilíbrios de subemprego" e racionalizar uma variedade de intervenções governamentais para "coordenar" movimentos em direção a equilíbrios superiores.

Mas a falha de coordenação sozinha não fornece um relato persuasivo das flutuações. Os modelos não têm papel para variáveis nominais e, na ausência de quaisquer barreiras ao ajuste de preços nominais, não há razão para que um distúrbio nominal afete a alocação real. Um choque nominal em um modelo de falha de coordenação não afeta o conjunto de alocações reais de equilíbrio. Assim, os modelos não fornecem uma base para a visão de que choques monetários e outros choques de demanda agregada são uma fonte crítica de flutuações na atividade econômica agregada.

De fato, se houver fricções no ajuste de preços, a questão de saber se existem múltiplos níveis de produto de equilíbrio ou simplesmente pequenos incentivos para as firmas ajustarem os preços na direção que move a economia em direção a um equilíbrio único de longo prazo é relativamente sem importância. Em qualquer caso, choques nominais têm efeitos reais, e esses efeitos provavelmente serão duradouros, uma vez que as forças que atuam para retornar o produto ao seu nível inicial são fracas.

Até mesmo a questão de se os choques têm efeitos permanentes não depende da questão de se existe um nível de equilíbrio único de produção. Na ausência de algum mecanismo através do qual os movimentos de produção de curto prazo alterem os níveis de produção de equilíbrio, os efeitos de qualquer distúrbio acabam desaparecendo, seja porque a economia retorna gradualmente ao nível único de produção de equilíbrio, ou porque ela é chocada aleatoriamente entre os diferentes níveis de equilíbrio. Na presença de tal mecanismo, por outro lado, choques nominais têm efeitos permanentes, independentemente de a produção ter um ou vários níveis de equilíbrio. Por exemplo, se o progresso tecnológico ocorre a partir do aprendizado pela prática (\textit{learning-by-doing}), o declínio na atividade causado por um choque nominal descendente e pelo ajuste lento de preços resultará em progresso tecnológico reduzido; o produto será, portanto, menor do que teria sido na ausência do choque, mesmo após os preços terem se ajustado totalmente. Modelos nos quais mudanças temporárias — aqui, movimentos da demanda agregada acoplados a barreiras temporárias ao ajuste de preços — têm efeitos permanentes são conhecidos como modelos de histerese (Blanchard e Summers, 1986).

\section*{O que são as Fricções?}

Obstáculos à flexibilidade nominal são centrais para os argumentos que tenho apresentado. Mas, além de observar que esses obstáculos são pequenos, disse pouco sobre a forma que poderiam assumir.\footnote{A discussão que se segue foca no ajuste imperfeito de preços nominais. Mas essencialmente os mesmos pontos se aplicam ao ajuste imperfeito de salários nominais, à informação imperfeita sobre o nível de preços ou a oferta de moeda, e à indexação incompleta de contratos de dívida.}

Começo com um ponto negativo: os custos de renegociar contratos, ou de coletar e processar informações e estimar o preço ótimo a cobrar, ou de informar clientes e fornecedores sobre um novo preço, não constituem, por si sós, custos de ajuste nominal de preços. A falácia subjacente à visão comum de que estes representam custos de mudança de preços nominais está em pensar que deixar seus preços nominais inalterados é a única maneira de uma firma "não fazer nada" em relação aos seus preços. Mas os preços poderiam ser ajustados por muitas regras simples — como um aumento de um determinado valor a cada mês ou a indexação a um índice de preços ou ao PIB nominal — que não envolveriam custos de renegociação, decisão ou comunicação uma vez implementadas. Para colocar de outra forma, o fato de que (por exemplo) é dispendioso informar os clientes sobre as mudanças nos preços nominais é uma consequência, não uma causa, do fato de que os preços nominais são geralmente deixados inalterados.

Mankiw foca nos "custos de menu" — os custos tecnológicos de mudar os preços nominais. (O exemplo padrão, é claro, é o custo incorrido por um restaurante ao imprimir novos menus — daí o nome.) Mas os custos de menu não podem explicar as evidências microeconômicas sobre a natureza das políticas de preços das firmas. O comportamento dos preços do catálogo da L. L. Bean, documentado por Kashyap (1991), ilustra as dificuldades. A Bean emite mais de 20 catálogos por ano, mas só altera os preços em dois catálogos (Outono e Primavera). Mesmo nesses catálogos, a maioria dos preços geralmente não é alterada. Nenhum dos fatos apoia a visão de que a barreira ao ajuste de preços é o custo de imprimir e postar um novo preço. Além disso, o tamanho das mudanças de preço varia tremendamente, e pequenas mudanças de preço têm a mesma probabilidade que grandes mudanças de serem seguidas rapidamente por uma mudança adicional. Finalmente, a frequência das mudanças de preço é baixa: em média, o preço de um bem só é alterado depois que a inflação corroeu o preço real em cerca de 10 por cento. Apenas um custo extremamente grande de ajuste de preços, ou um custo extremamente pequeno de não cobrar o preço que seria ótimo na ausência de custos de ajuste, pode reconciliar esta descoberta com uma visão de "custo de menu". Outros estudos da microeconomia do ajuste de preços relatam descobertas semelhantes (Carlton, 1986; Cecchetti, 1986).

Akerlof e Yellen descrevem a barreira ao ajuste nominal como "quase-racionalidade". Ou seja, eles sugerem que as firmas simplesmente abrem mão de pequenas quantias de lucros. Mas muitas políticas de preços — na verdade infinitas — envolvem pequenas quantias de lucros perdidos. A questão é por que, de todas essas políticas, as firmas escolhem aquelas que envolvem rigidez nominal considerável; como Akerlof e Yellen observam, a observação de que a rigidez nominal tem apenas pequenos custos não ajuda a responder a essa pergunta. No entanto, a sugestão de que a fricção pode assumir a forma de "quase-racionalidade" é importante: ela sugere que as barreiras à flexibilidade nominal não precisam ser puramente tecnológicas.

O que atualmente parece ser o caminho mais promissor parte da observação de McCallum (1986) de que, como bens e trabalho são geralmente trocados por dólares, e não por outros bens, é computacionalmente mais fácil postar preços e salários em unidades monetárias. Em outras palavras, é natural usar o meio de troca como unidade de conta. Assim, os preços são postados em termos nominais em vez de reais. Os custos de menu — os custos físicos de mudar os preços — fazem com que os preços não sejam ajustados continuamente. A simplicidade computacional e os custos de menu, por si sós, dificilmente gerariam uma rigidez nominal substancial. Mas eles podem ter o efeito de fazer com que manter os preços nominais inalterados se torne equivalente a "não fazer nada" e, assim, gerar uma rigidez considerável.

Em particular, se os preços são normalmente deixados fixos, ajustar um preço para levar em conta movimentos da demanda agregada (seja através de mudanças de preço ou através da adoção de algum tipo de mecanismo de indexação) requer uma decisão consciente por parte do definidor de preços. As barreiras à flexibilidade total de preços incluem, então, não apenas custos de computação e de menu, mas também a necessidade de os definidores de preços perceberem os benefícios de ajustar seus preços em resposta a choques de demanda. Além disso, se a maioria das firmas ajusta seus preços apenas infrequente, os custos para uma única firma ao adotar uma política de preços diferente incluem não apenas os custos diretos, mas também os custos de explicar aos clientes qual é a política de preços e como ela opera. (Imagine, por exemplo, a L. L. Bean incluindo em seus catálogos um aviso de que todos os preços são indexados ao Índice de Preços ao Consumidor, ou aumentados automaticamente em um centésimo de um por cento por dia.) Finalmente, se os preços são postados e os salários são pagos em dólares e não são ajustados continuamente, os indivíduos podem vir a atribuir importância aos preços e salários em dólares — um preço ou salário nominal inalterado pode vir a ser visto como a norma.

O resultado final é provável que custos que em um mundo sem fricções estariam associados a mudanças nos preços reais — custos de coletar e processar informações e chegar a uma decisão, custos de negociação, custos de ofender clientes e empregados que preferem preços e salários estáveis, e assim por diante — tornem-se vinculados, em parte, a mudanças nos preços nominais. A rigidez nominal pode, portanto, ser mais forte e mais complicada do que seria se a computação e os custos de menu fossem as únicas fricções.

Esta análise sugere que a taxa de inflação é um determinante importante da força das fricções no ajuste de preços. Se a inflação é alta, os preços nominais são ajustados com frequência, os definidores de preços aprendem que é importante ajustar seus preços aos movimentos do nível de preços agregado, e os indivíduos passam a não atribuir grande importância aos preços e salários nominais. As teorias de rigidez de preços decorrentes de pequenas fricções nominais, portanto, preveem que os efeitos reais de um dado choque nominal são menores em ambientes com inflação predominante mais alta. Esta previsão difere das previsões de teorias concorrentes. Nas teorias keynesianas tradicionais, o grau de rigidez nominal é exógeno; na teoria da informação imperfeita de Lucas (1973), o grau de rigidez nominal é determinado pelas variâncias dos choques de demanda específicos da firma e agregados e não depende do nível de inflação. Ball, Mankiw e Romer (1988) examinam os efeitos reais dos movimentos da demanda agregada tanto entre países quanto dentro de países em diferentes períodos de tempo; eles descobrem que, como preveem as teorias novo-keynesianas, os efeitos reais dos deslocamentos de demanda são menores em ambientes mais inflacionários.

\section*{Direções de Pesquisa}

O argumento novo-clássico de que a suposição de rigidez nominal na macroeconomia keynesiana era teoricamente incoerente — não apenas que faltavam fundamentos microeconômicos, mas que era inconsistente com quaisquer suposições microeconômicas defensáveis — foi refutado. Neste ponto, a questão relevante não é mais se tais fundamentos podem ser fornecidos, mas se os modelos construídos sobre esses fundamentos descrevem corretamente a realidade. A pesquisa direcionada a responder a essa pergunta está prosseguindo em três frentes.

Uma primeira área de pesquisa envolve as próprias fricções. Como sugere a seção anterior, houve apenas alguns estudos de ajuste de preços ao nível de firmas individuais, e seus resultados são intrigantes. Além disso, as firmas e bens estudados são idiossincráticos, e não sabemos se estudos de outras firmas produziriam descobertas semelhantes. Em suma, os economistas não têm uma boa compreensão das políticas de ajuste de preços das firmas, ou mesmo das considerações que fundamentam suas escolhas de políticas.

A natureza exata das fricções microeconômicas no ajuste de preços provavelmente terá implicações importantes para a rigidez nominal ao nível macroeconômico. Caplin e Spulber (1987) apresentam um exemplo em que os custos microeconômicos de ajuste de preços não levam a nenhuma rigidez nominal no agregado. Suas suposições principais são que todas as mudanças de preço são aumentos e que qualquer firma aumenta seu preço sempre que a demanda agregada aumentou em um valor fixo desde seu aumento de preço anterior. Eles demonstram que, com essas suposições, é possível que a fração de firmas fazendo ajustes de preços em qualquer momento varie com o tamanho dos choques nominais exatamente da maneira necessária para fazer com que os choques deixem o produto real global inalterado.

Outras suposições sobre as características microeconômicas das fricções, no entanto, têm implicações muito diferentes para o comportamento macroeconômico. Por exemplo, as políticas de ajuste das firmas podem conter um elemento importante de tempo fixo; isto é, até certo ponto, uma determinada firma pode mudar seu preço após um período fixo de tempo, em vez de após a demanda agregada ter mudado por um valor fixo. Se tais políticas de preços forem acopladas a uma rigidez real considerável, os distúrbios nominais terão efeitos reais que não são apenas grandes, mas também duradouros. Neste caso, a fração de firmas mudando os preços responde pouco a um choque, e as rigidezes reais fazem com que as firmas que se ajustam façam apenas pequenas mudanças de preço (por exemplo, Blanchard, 1983). Dado o quão criticamente as consequências macroeconômicas dos choques nominais dependem das especificidades das políticas de ajuste de preços das firmas, os estudos sobre as barreiras ao ajuste de preços são claramente um assunto premente para a pesquisa.

A segunda área de pesquisa são as rigidezes reais. A existência de fatores que fazem com que as firmas desejem apenas pequenas mudanças em seus preços relativos em resposta aos movimentos da demanda agregada na produção é necessária se os choques de demanda tiverem efeitos reais substanciais, em vez de efeitos principalmente sobre os preços; essas rigidezes reais provavelmente devem incluir razões para grandes movimentos no emprego e horas serem acompanhados por apenas pequenas mudanças nos salários reais. Salários de eficiência, externalidades de mercado denso, imperfeições no mercado de capitais, o comportamento cíclico das elasticidades da demanda e uma variedade de outras fontes potenciais de rigidez real são temas ativos de pesquisa teórica e empírica.

Estudar fenômenos microeconômicos como fricções e rigidezes reais, no entanto, dificilmente será suficiente. Se estudarmos fenômenos microeconômicos sem atenção aos fenômenos macroeconômicos que estamos tentando entender, poderemos razoavelmente concluir que as pequenas fricções no ajuste nominal não são importantes. Mas quando nos voltamos para a macroeconomia, ficamos intrigados com os grandes efeitos reais dos choques de demanda agregada. Compreender como as propriedades microeconômicas da economia dão origem aos fenômenos macroeconômicos observados é um objetivo realista. Mas unir a microeconomia e a macroeconomia pode não ser possível: as simplificações que são úteis para entender a maioria dos fenômenos microeconômicos podem ser fatais para os esforços de entender as flutuações macroeconômicas.

Portanto, um terceiro caminho de pesquisa, talvez o mais óbvio e fundamental, é examinar as evidências macroeconômicas relativas aos efeitos da moeda e de outros distúrbios da demanda agregada. Estudos sobre se os movimentos da demanda agregada são centrais para as flutuações reais têm sido um foco principal da pesquisa macroeconômica. Dada a sua dificuldade e a importância da questão que abordam, eles certamente continuarão a sê-lo.

\section*{Conclusão}

As escolas de pensamento sobre as flutuações macroeconômicas podem ser classificadas de forma útil de acordo com suas respostas a duas perguntas. A primeira é se a dicotomia clássica falha; a segunda é se a economia possui características não walrasianas que são importantes para as flutuações além das suposições de falha da dicotomia clássica. A classificação resultante dois-por-dois é mostrada na Figura 3.

Praticamente toda a macroeconomia de meados da década de 1970 se encaixa no quadro inferior esquerdo do diagrama. Nas teorias keynesianas e monetaristas, nos modelos de desequilíbrio e até mesmo na teoria da informação imperfeita de Lucas, o único desvio importante das suposições walrasianas é a presença de alguma imperfeição — geralmente o ajuste lento de preços ou salários nominais, ou informação imperfeita sobre distúrbios reais e nominais — que faz com que distúrbios nominais tenham efeitos reais.

Os três quadros restantes do diagrama mostram as três principais escolas atuais de pensamento sobre flutuações econômicas. Todas as três escolas discordam da visão incontestada de 15 anos atrás em relação à resposta de pelo menos uma das duas perguntas consideradas no diagrama. As teorias dos ciclos reais de negócios — que se baseiam nas premissas de que a economia é essencialmente walrasiana e que a dicotomia clássica se mantém — são mostradas no canto superior esquerdo do diagrama. O que chamo de modelos de "falha de coordenação" está no canto superior direito. Incluo nesta escola tanto os modelos de múltiplos equilíbrios — por exemplo, o modelo de Diamond de externalidades de mercado denso — quanto modelos nos quais as imperfeições reais são centrais para as flutuações, mas existe um equilíbrio único. Esta escola atribui um papel central às características não walrasianas da economia, mas nenhum aos distúrbios nominais. Finalmente, os modelos que foram o tema principal deste artigo, que repousam na crença de que tanto a falha da dicotomia clássica quanto os elementos não walrasianos da economia são essenciais para o ciclo de negócios, estão no quadro inferior direito do diagrama.

\begin{center}
\textbf{Figura 3} \\
\textbf{Escolas de Pensamento Macroeconômico} \\
\begin{tabular}{|l|c|c|}
\hline
& \multicolumn{2}{c|}{\textbf{A economia possui características}} \\
& \multicolumn{2}{c|}{\textbf{não walrasianas importantes?}} \\
\cline{2-3}
& \textbf{Não} & \textbf{Sim} \\
\hline
\textbf{A dicotomia clássica falha?} & & \\
\textbf{Não} & Teorias dos Ciclos & Teorias de Falha de \\
& Reais de Negócios & Coordenação \\
\hline
\textbf{Sim} & Teorias Tradicionais Keynesianas & Teorias \\
& e Monetaristas; Teoria da Informação & Novo-Keynesianas \\
& Imperfeita de Lucas & \\
\hline
\end{tabular}
\end{center}

Das três áreas ativas de trabalho mostradas na Figura 3, apenas os modelos novo-keynesianos fornecem uma explicação para a importância dos distúrbios nominais para a economia real; e, como sugeri no início, eles também fornecem a explicação mais plausível de por que outros choques de demanda agregada importam. Assim, a menos que novos trabalhos empíricos derrubem a visão amplamente compartilhada da importância da moeda e de outros choques de demanda agregada, a análise descrita neste artigo será uma parte necessária de qualquer modelo completo de flutuações macroeconômicas.

\begin{itemize}
\item \textit{Agradeço a George Akerlof, Laurence Ball, N. Gregory Mankiw, Christina Romer, Carl Shapiro, Joseph Stiglitz, Timothy Taylor e Janet Yellen pelos comentários úteis, e à National Science Foundation e à Fundação Sloan pelo apoio financeiro.}
\end{itemize}

\section*{Referências}

\begin{description}
\item Akerlof, George A., e Janet L. Yellen, "A Near-Rational Model of the Business Cycle, with Wage and Price Inertia," \textit{Quarterly Journal of Economics}, Suplemento, 1985, 100:5, 823–38.
\item Altonji, Joseph G., "Intertemporal Substitution in Labor Supply: Evidence from Panel Data," \textit{Journal of Political Economy}, Junho 1986, Parte 2, 94, S176–S215.
\item Ball, Laurence, N. Gregory Mankiw, e David Romer, "The New Keynesian Economics and the Output-Inflation Trade-off," \textit{Brookings Papers on Economic Activity}, 1988, no. 1, 1–65.
\item Ball, Laurence, e David Romer, "Are Prices Too Sticky?" \textit{Quarterly Journal of Economics}, Agosto 1989, 104, 507–24.
\item Ball, Laurence, e David Romer, "Real Rigidities and the Non-Neutrality of Money," \textit{Review of Economic Studies}, Abril 1990, 57, 183–203.
\item Bernanke, Ben, e Mark Gertler, "Agency Costs, Net Worth, and Business Fluctuations," \textit{American Economic Review}, Março 1989, 79, 14–31.
\item Blanchard, Olivier J., "Price Asynchronization and Price Level Inertia." Em Dornbusch, Rudiger, e Mario Henrique Simonsen, eds., \textit{Inflation, Debt, and Indexation}. Cambridge: MIT Press, 1983, 3–24.
\item Blanchard, Olivier J., e Lawrence H. Summers, "Hysteresis and the European Unemployment Problem," \textit{NBER Macroeconomics Annual}, 1986, 1, 15–78.
\item Caplin, Andrew S., e Daniel F. Spulber, "Menu Costs and the Neutrality of Money," \textit{Quarterly Journal of Economics}, Novembro 1987, 102, 703–25.
\item Carlton, Dennis W., "The Rigidity of Prices," \textit{American Economic Review}, Setembro 1986, 76, 637–58.
\item Cecchetti, Stephen G., "The Frequency of Price Adjustment: A Study of the Newsstand Prices of Magazines," \textit{Journal of Econometrics}, Agosto 1986, 31, 255–74.
\item Cooper, Russell, e Andrew John, "Coordinating Coordination Failures in Keynesian Models," \textit{Quarterly Journal of Economics}, Agosto 1988, 103, 441–63.
\item Diamond, Peter A., "Aggregate Demand Management in Search Equilibrium," \textit{Journal of Political Economy}, Outubro 1982, 90, 881–94.
\item Kashyap, Anil K., "Sticky Prices: New Evidence from Retail Catalogs," manuscrito não publicado, University of Chicago Graduate School of Business, Dezembro 1991.
\item Lucas, Robert E., Jr., "Some International Evidence on Output-Inflation Tradeoffs," \textit{American Economic Review}, Junho 1973, 63, 326–34.
\item Malinvaud, Edmond, \textit{The Theory of Unemployment Reconsidered}. Oxford: Basil Blackwell, 1977.
\item Mankiw, N. Gregory, "Small Menu Costs and Large Business Cycles: A Macroeconomic Model of Monopoly," \textit{Quarterly Journal of Economics}, Maio 1985, 100, 529–37.
\item McCallum, Bennett T., "On 'Real' and 'Sticky-Price' Theories of the Business Cycle," \textit{Journal of Money, Credit, and Banking}, Novembro 1986, 18, 397–414.
\item Meltzer, Allan H., \textit{Keynes's Monetary Theory: A Different Interpretation}. Cambridge: Cambridge University Press, 1988.
\item Plosser, Charles I., "Understanding Real Business Cycles," \textit{Journal of Economic Perspectives}, Verão 1989, 3:3, 51–77.
\item Reaume, David M., "On the Misuse of Taylor's Theorem in Economics," manuscrito não publicado, University of Alaska, Southeast, 1991.
\end{description}

\end{document}